\subsection*{Zero-padding: \texttt{.zfill}}

Every tiled kernel faces the same edge case: what happens when M or K is not a multiple of the tile size? The last tile along that axis is partial --- some of its elements lie beyond the tensor boundary.

A separate code path for the boundary tile adds branches and complicates a uniform MMA loop. The \texttt{.zfill} modifier handles the missing source elements during transfer:

\begin{lstlisting}[language=croqtile,caption={Zero-padding: .zfill}]
tma.copy.swiz<3> lhs.subspan(WARP_M, TILE_K).at(block_m, iv_k).zfill
  => lhs_load_s;
\end{lstlisting}

The \co@.zfill@ modifier supplies zero values where the destination tile extends beyond the valid source region. In a matrix product, those zeros add no contribution to the sum. The consumer can therefore process a full tile with a known value in every position. This establishes the required numerical behavior; the cost of handling the boundary still depends on the target and transfer path.

\begin{figure}[htbp]\centering
\includegraphics[width=\linewidth,height=0.56\textheight,keepaspectratio]{f09-fig6_zfill_light.png}
\caption{A [4, 4] boundary tile from a [6, 8] source keeps its valid rows and fills the two out-of-range rows with zeros.}
\end{figure}
